% fig_sqw.tex -- native pgfplots (figure*, full width). Charge S(q,w) + spin S^zz(q,w) heatmaps.
% L=12 sector-Lanczos/Haydock (spectral_lanczos.py). Data field = rasterized PNG (paperhot + gamma);
% axes/colorbar/guides = native. Each panel on its OWN frequency window (charge ~U up to 10t; spin ~J up
% to 2.4t) so neither channel is crushed -- the different windows ARE the spin-charge-separation signature.
% 2026-09-27 (B9): the PNG rasters span the FULL computed grids (charge omega in [0, 11]t, spin [0, 2.6]t, q/pi in
% [qmin, qmax]; src/export_sqw_dat.py -> sqw_extent.dat). They used to be placed at ymax=10 / 2.4, which drew every
% feature at 10/11 = 0.909x and 2.4/2.6 = 0.923x of its true frequency. The extents are now read from
% sqw_extent.dat; the axis windows (ymax=10, 2.4) are unchanged and clip the excess.
\pgfplotstableread{figs/sqw_extent.dat}\sqwextent
\pgfplotstablegetelem{0}{qmin}\of\sqwextent\edef\sqwQmin{\pgfplotsretval}
\pgfplotstablegetelem{0}{qmax}\of\sqwextent\edef\sqwQmax{\pgfplotsretval}
\pgfplotstablegetelem{0}{wc_max}\of\sqwextent\edef\sqwWCmax{\pgfplotsretval}
\pgfplotstablegetelem{0}{ws_max}\of\sqwextent\edef\sqwWSmax{\pgfplotsretval}
\begin{tikzpicture}
% ---- charge S(q,omega): gapped, weight near omega ~ U ; the empty low-omega region IS the charge gap ----
\begin{axis}[
  name=charge, paperaxis,
  width=7.7cm, height=6.6cm,
  grid=none, enlargelimits=false, axis on top,
  xmin=0.1667, xmax=1.8333, ymin=0, ymax=10,
  ytick={0,2,4,6,8,10}, xtick={0.3333,0.6667,1,1.3333,1.6667},
  xticklabels={$\tfrac13$,$\tfrac23$,$1$,$\tfrac43$,$\tfrac53$},
  xlabel={momentum transfer $q/\pi$}, ylabel={frequency $\omega/t$},
  title={\small charge $S(q,\omega)$ \, (gapped, scale $\sim U$)},
  colormap name=paperhot, colorbar, point meta min=0, point meta max=1,
  colorbar style={width=0.22cm, ytick={0,0.5,1}, title={\scriptsize norm.}, tick label style={font=\scriptsize}},
]
  \addplot graphics[xmin=\sqwQmin, xmax=\sqwQmax, ymin=0, ymax=\sqwWCmax] {figs/sqw_charge.png};
  % the charge gap = the forbidden band from omega=0 up to the onset Delta_c: brace it and name it
  \addplot[white, opacity=0.9, densely dashed, line width=0.7pt, forget plot] coordinates {(0.1667,5.71) (1.8333,5.71)};
  \draw[decorate, decoration={brace, amplitude=6pt}, white, line width=0.9pt]
      (axis cs:0.30,5.63) -- (axis cs:0.30,0.06);
  \node[white, font=\scriptsize, anchor=west, align=left] at (axis cs:0.37,2.7)
      {charge gap\\ $\Delta_c\!\approx\!5.7\,t$\\ (no weight)};
  \node[white, font=\scriptsize\itshape, anchor=west] at (axis cs:0.22,6.05) {onset};
\end{axis}
% ---- spin S^zz(q,omega): gapless spinon continuum, scale ~ J, fills the panel ----
\begin{axis}[
  name=spin, paperaxis,
  at={(charge.east)}, anchor=west, xshift=2.35cm,
  width=7.7cm, height=6.6cm,
  grid=none, enlargelimits=false, axis on top,
  xmin=0.1667, xmax=1.8333, ymin=0, ymax=2.4,
  ytick={0,0.5,1,1.5,2}, xtick={0.3333,0.6667,1,1.3333,1.6667},
  xticklabels={$\tfrac13$,$\tfrac23$,$1$,$\tfrac43$,$\tfrac53$},
  xlabel={momentum transfer $q/\pi$}, ylabel={frequency $\omega/t$},
  title={\small spin $S^{zz}(q,\omega)$ \, (gapless, scale $\sim J$)},
  colormap name=paperhot, colorbar, point meta min=0, point meta max=1,
  colorbar style={width=0.22cm, ytick={0,0.5,1}, title={\scriptsize norm.}, tick label style={font=\scriptsize}},
]
  \addplot graphics[xmin=\sqwQmin, xmax=\sqwQmax, ymin=0, ymax=\sqwWSmax] {figs/sqw_spin.png};
  % des Cloizeaux--Pearson two-spinon edges -- thin DOTTED guide (Paper-2 style), subtle not dominant
  \addplot[white, opacity=0.8, densely dotted, line width=0.7pt, forget plot] table[x=q,y=om] {figs/dcp_lower.dat};
  \addplot[white, opacity=0.8, densely dotted, line width=0.7pt, forget plot] table[x=q,y=om] {figs/dcp_upper.dat};
  % extracted dominant spinon peak per q: data tracking the dCP dispersion
  \addplot[only marks, mark=*, mark size=1.3pt, forget plot,
           mark options={fill=pTeal, draw=pCream, line width=0.4pt}] table[x=q,y=om] {figs/sqw_spinpeak.dat};
  \node[white, font=\scriptsize\itshape, anchor=west, align=left] at (axis cs:0.20,2.18)
      {gapless spinon continuum};
  \node[white, font=\scriptsize, anchor=east] at (axis cs:1.82,1.98) {dCP guide (dotted) {\normalfont\&} extracted peaks};
\end{axis}
\end{tikzpicture}
